\section{Introduction}

\subsection{The Big Red Button as a Cultural Interface}

A big red button is rarely allowed to remain a button for very long. In ordinary interface terms it is a surface that can be pressed, but in cultural terms it is a compressed scene of decision, danger, permission, and theatrical consequence. Button histories show that the push-button became powerful partly because it made action appear simple while relocating technical complexity elsewhere \citep{Plotnick2018}. That history matters for the present study because the big red button is not merely a large target in virtual space. It is a familiar interface form whose apparent simplicity has already been trained by domestic electrification, industrial control, emergency equipment, software widgets, political metaphor, and popular comedy.

The red button is especially ambivalent because it carries two opposed promises. On one side, it is the button of doom: the nuclear launch button, the alarm, the forbidden switch, the small object that makes catastrophe imaginable as a single gesture. Push-button warfare has long been represented as distant action condensed into remote control \citep{PlotnickWarfare2012}, and nuclear discourse often compresses diffuse institutions of command, strategy, and threat into compact artifacts such as buttons, clocks, and briefcases \citep{Taylor2007NuclearDiscourse,Cheek2025ApocalypticTechnicalCommunication}. On the other side, the same visual grammar appears in emergency-stop design, where a large, salient, reachable control is meant to prevent catastrophe rather than initiate it \citep{Kwon1996EmergencyStop}. The button therefore does not simply mean danger. It also means intervention. It can be the problem, the solution, and the conveniently circular object on which both interpretations have agreed to sit.

This ambivalence gives the big red button its peculiar narrative force. In fiction, cartoons, memes, and interface jokes, a red button often functions like a materialized foreshadowing device: once it appears, the scene has already begun asking whether it will be pressed. The familiar ``Daily Struggle'' or ``Two Buttons'' image macro condenses this structure into a cartoon figure visibly sweating between two red-button options \citep{DailyStruggleKnowYourMeme}. Meme scholarship helps explain why that image works: internet memes are not isolated jokes but repeatable cultural units that invite imitation, variation, and participation \citep{Shifman2013Memes}, while memetic media circulate through multimodality, reappropriation, resonance, collectivism, and spread \citep{Milner2016WorldMadeMeme}. The two-buttons meme is thus not just a picture of indecision. It is a small cultural machine for making conflicted action visible, legible, and shareable.

The present paper treats this cultural saturation as an empirical problem rather than as decoration. A participant who encounters a large red button in a neutral virtual room is not encountering a neutral object that happens to be red. The object arrives already mediated by histories of domestic convenience, emergency control, nuclear rhetoric, platform reaction buttons, meme templates, and comedic prohibition. This does not mean every participant will interpret the button in the same way, and the paper explicitly avoids treating red-button symbolism as universal. It does mean that voluntary pressing in this paradigm must be analyzed as a multilevel relation among affordance, cultural memory, bodily reach, curiosity, agency, humor, and physiological coupling. We call that relation \emph{pressability}.

\subsection{Button Histories and the Delegation of Action}

The modern push-button should be understood as a late nineteenth-century interface for electrification rather than as a timeless control primitive. Br{\aa}then and Herstad describe the electric push-button as a switching device that emerged in the last two decades of the nineteenth century, when push-buttons enabled immediate control of electrical devices \citep{BrathenHerstad2023Pushbutton}. Plotnick shows that between 1880 and 1923 push-buttons spread rapidly through domestic and public electrical systems, including bells, lights, elevators, and service devices \citep{PlotnickInterface2012}. This period matters because the button moved control from exposed mechanism or specialist operation to a small standardized surface. The user no longer needed to see or understand the whole system in order to act on it. The finger became an acceptable delegate.

Early push-buttons also changed what users were expected to know. Plotnick argues that from roughly 1880 to 1915 push-buttons were often framed as part of electrical education, since they introduced lay users to the possibilities and risks of domestic electrical systems \citep{PlotnickInterface2012}. After about 1915, growing familiarity helped stabilize the button as a black-box interface, meaning that the user could treat the button as a practical endpoint rather than as an invitation to understand the circuit behind it \citep{PlotnickInterface2012}. The historical shift is important for HCI because it marks an early form of user-facing abstraction. Complexity was not removed. It was made someone else's immediate problem, which is one of interface design's most enduring courtesies.

The push-button also changed who could appear to control a machine. Plotnick's broader history of button pushing shows that buttons were made, distributed, celebrated, rejected, and refashioned across social settings \citep{Plotnick2018}. Work on haptics and buttons further shows that force, flatness, and tactile feedback shape how users understand action through machines \citep{Plotnick2017HapticsButtons}. Janlert describes the button as a ubiquitous interaction form, one whose ordinary visibility makes it easy to overlook as a theoretical object \citep{Janlert2014UbiquitousButton}. Duffner's organizational account of the push-button frames it as a mediating technology, which means that the press does not simply operate a machine but helps organize relations among people, work, and technical systems \citep{Duffner2019PushButtonMediating}. A button is therefore never only a switch. It is also a small organizational settlement about who may act, how quickly, and with how little explanation.

The transition from mechanical to digital buttons intensifies this settlement. Mechanical buttons often close or break circuits, and their state can be visible or tactile \citep{BrathenHerstad2023Pushbutton}. Digital buttons, by contrast, become inputs into programmed systems, so the same visible or tactile gesture can trigger processing that is conditional, delayed, distributed, or opaque \citep{BrathenHerstad2023Pushbutton}. HCI inherited this useful ambiguity by making computation appear graspable through visible and actionable objects. Shneiderman's account of direct manipulation emphasizes visible objects, rapid feedback, incremental action, and reversibility \citep{Shneiderman1983DirectManipulation}, while Myers traces how graphical interfaces stabilized widgets such as buttons, menus, and scrollbars as ordinary software components \citep{Myers1998HistoryHCI}. The continuity of the press hides a discontinuity in control: the hand performs an old gesture, while software decides what that gesture now means.

\subsection{Catastrophe, Shutdown, and Ambivalent Command}

The big red button is the push-button after promotion. It intensifies ordinary pressability into a public scene of possible consequence, where the visible action is small and the imagined effect is large. Nuclear writing and reactor writing both use big-button imagery because it condenses control into a legible form \citep{Tucker2019BigRedButton}. Political communication around nuclear command similarly depends on artifacts that stand in for institutions too distributed to picture easily \citep{Taylor2007NuclearDiscourse}. The result is a cultural shorthand in which a single press appears to summarize command, authority, risk, and temperament. This shorthand is analytically useful precisely because it is false in the literal sense. Nuclear command is not a single magic button, but public imagination repeatedly prefers the button because distributed procedure is less photogenic.

The button's catastrophic associations are not confined to formal politics. Popular media and games repeatedly turn nuclear action into symbolic, playable, or spectacular form. Pantoliano argues that video games can convert nuclear weapons into rewards, success conditions, and objects of pleasurable mastery, while downplaying humanitarian harm \citep{Pantoliano2025Playing}. Faux and Pullen likewise show that contemporary nuclear narratives can become occasions for rethinking how nuclear risk is represented in popular culture \citep{FauxPullen2025Catalyst}. In this broader media ecology, the big red button becomes a device for making the unthinkable operational. It lets culture imagine catastrophe as a user interface problem, which is an alarming convenience and therefore a useful object for research.

The same form also belongs to safety. Emergency-stop buttons are designed to be visually salient, reachable, and rapidly actionable under stress \citep{Kwon1996EmergencyStop}. Automotive and industrial control research similarly treats emergency button layout as a matter of human performance under time pressure \citep{ChoiEtAl2010EmergencyButtonLayout}. Recent writing on nuclear safety even imagines an ``undo button'' for nuclear armageddon as a technical and policy intervention into irreversible command \citep{AllenLevin2025UndoButton}. The red button therefore oscillates between initiation and interruption. It can launch the disaster or stop the machine. The interface is the same shape because both scenarios ask the body to act before deliberation has finished drafting its minutes.

This double role matters for the VR study because the participant is not simply deciding whether to activate a stimulus. They are encountering an object that has inherited the cultural logic of both forbidden action and emergency rescue. The ambiguity is not a flaw in the paradigm. It is the paradigm. A button that only meant one thing would be easier to model and less interesting to press. A big red button means too much, and the experiment asks what participants do when that excess of meaning is made reachable, countable, and physiologically responsive.

\subsection{Media, Self, and Physiological Extension}

The study's embodied-feedback condition draws on a longer media-theoretical problem: technologies do not only serve users, but reorganize the relation between bodies, senses, and environments. McLuhan's account of media as extensions of human faculties gives this paper a useful starting point, since a medium matters not only because of its content but because it changes the scale, pace, and pattern of human action \citep{McLuhan1964UnderstandingMedia}. A button is modest compared with print, radio, or television, but it is also a medium in McLuhan's sense: it extends the finger into a system that acts beyond the finger's physical reach. The red button is therefore not simply an object in the scene. It is a hinge between bodily action and delegated consequence.

The concept of the extended self sharpens this point by showing that persons can experience external objects, possessions, and digital traces as implicated in selfhood. Belk's classic account of the extended self argues that possessions can become part of how people understand and enact identity \citep{Belk1988ExtendedSelf}. Belk later revises that account for digital environments, where self-extension occurs through dematerialized, networked, and datafied forms \citep{Belk2013ExtendedSelfDigital}. The BRB study does not assume that a virtual button becomes part of the self in this strong sense. Instead, it asks a narrower and more testable question: can a button become more self-adjacent when it visibly pulses with live cardiac timing and when the surrounding environment brightens and dims with respiration?

Virtual reality makes that question experimentally tractable because VR can reorganize the perceived relation between body, action, and environment. Kilteni, Groten, and Slater define embodiment in VR through body ownership, agency, and self-location \citep{KilteniEtAl2012}. Serino describes peripersonal space as a multisensory interface between the individual and the environment, where perception and possible action are coupled in the near-body zone \citep{Serino2019}. VR research on peripersonal space extends this account to real, virtual, and mixed environments \citep{SerinoEtAl2018}. If a reachable virtual object responds to physiological signals, then the object may shift from being merely visible to being body-adjacent in the participant's action field. That shift is not mystical. It is the sober possibility that the button has learned to keep time with the user.

The implementation therefore uses live and sham physiological feedback as a test of embodied coupling rather than as decorative biometric theater. One condition maps Polar H10 electrocardiographic timing to a red pulse or tone associated with the button, while a control condition uses simulated ECG-like activity generated through NeuroKit2-style signal synthesis \citep{MakowskiEtAl2021}. Respiration-derived feedback modulates ambient lighting, with inhalation brightening the room and exhalation dimming it. VR biofeedback research shows that physiological feedback can be integrated into immersive environments \citep{LueddeckeFelnhofer2022}, heart-rate-variability biofeedback has been implemented in VR \citep{RockstrohEtAl2019}, and real versus placebo breathing feedback can produce different experiential outcomes in VR relaxation contexts \citep{ChittaroEtAl2024}. The BRB adaptation is intentionally stranger: it asks whether physiology can increase the pressability of an object whose cultural affordance is already loud enough to file a noise complaint.

\subsection{Pressability, Curiosity, and Social Legibility}

Pressability names the relation produced when an object is perceived as available for pressing and when that availability matters. Affordance theory gives HCI a language for action possibilities \citep{Gaver1991}, and Norman's account of perceived affordance connects artifact cues to expected user action \citep{Norman1999}. Experience-centered HCI adds that action is interpreted through felt situation and meaning rather than through mechanical input alone \citep{McCarthyWright2004}. In this paper, pressability therefore refers to more than target size, color, or reach. It refers to the multilevel condition under which a person experiences the button as something that can, might, should, should not, or absolutely must not be pressed.

Voluntary pressing under minimal task demand also depends on motivation without instruction. Curiosity research treats information gaps and uncertainty as drivers of exploratory behavior \citep{KiddHayden2015}. Boredom and under-stimulation can push people toward self-generated events even when those events are not useful in any narrow instrumental sense \citep{Wilson2014}. The BRB environment deliberately gives participants a sparse room, a single highly salient object, and explicit permission to press or abstain. This preserves non-pressing as a valid outcome while making pressing available as a way to resolve uncertainty, create an event, or participate in the situation. The button does not command. It waits, which is worse.

The visible counter above the button makes pressing socially and methodologically legible. Sense of agency concerns the relation between action and experienced consequence \citep{HaggardChambon2012}; in the BRB paradigm, the counter gives each press an immediate trace, so the participant can see action becoming data. Social facilitation research in HCI shows that performance can change when interaction is socially framed or when a system appears responsive \citep{HallHenningsen2008}. Digital reaction-button research likewise shows that button presses can translate affect into countable platform signals \citep{GeboersEtAl2020}. The counter therefore does not merely report successful input. It changes the press from a private motor act into a visible contribution to an accumulating record. The participant becomes a co-producer of the dataset, which is a dignified way to say that the button has acquired paperwork.

The button's placement and color complete the action invitation. The current VR prototype positions the button approximately 15 degrees below eye level, roughly 0.75\,m from the viewpoint, and at about 10\,cm diameter, grounding the object in prior work on head-worn display placement, visual layout, and arm's-length reach \citep{Lin2021,Alghofaili2019,Bonin2022}. Red can operate as an arousal, warning, or avoidance cue, although color meanings vary across cultures and should not be universalized \citep{ElliotMaier2014,GaoEtAl2007}. Feature-singleton research shows that visual contrast can capture attention before deliberate interpretation is complete \citep{Li2008}, and work on information-optimal local features similarly supports the attentional force of salient local contrast \citep{CastellottiEtAl2022}. The BRB is therefore designed as a saturated red singleton against a low-contrast background. It is visually polite only in the sense that a fire alarm is polite.

\subsection{Humor as an Interface to Serious Interaction}

Humor is not a detour from the paper's interface argument, because humor exposes what interfaces usually try to stabilize: the relation between expectation and consequence. Minsky's cognitive account of jokes treats humor as a way of encountering and managing otherwise troublesome or defective patterns of thought \citep{Minsky1980Jokes}. Contemporary computational humor research similarly treats humor as difficult because it depends on situation, timing, interpretation, and shared assumptions rather than on wording alone \citep{Cowie2023ComputationalHumor,LemmensDeMarez2026HumorSurvey}. A joke and an interface both depend on what the user expects to happen next. Humor becomes relevant to HCI when the system uses expectation violation as part of interaction rather than as a system failure.

HCI has repeatedly taken humor seriously, with the expression appropriate to the field. Morkes, Kernal, and Nass directly tested humor in task-oriented human-computer interaction and computer-mediated communication \citep{MorkesEtAl1999HumorHCI}. Dybala and colleagues developed humorized computational intelligence for user-adapted systems \citep{DybalaEtAl2009HumorizedCI}. Nijholt treats humor as a movement from word play to world play in interaction with computational agents, robots, and environments \citep{Nijholt2018WordPlayWorldPlay}. Conversational-agent studies also show both the promise and risk of machine humor, since humor can support engagement but can fail when systems misread social context or overclaim rapport \citep{ZarghamEtAl2023FunnyHow,SunEtAl2024FunnyChatbot}. Humor is therefore an interface modality in the modest sense: it shapes timing, expectation, affect, and willingness to continue interacting.

The BRB also belongs to design lineages that treat uselessness, ambiguity, and absurdity as resources for inquiry. Gaver, Beaver, and Benford argue that ambiguity can be a resource for design because it invites interpretation rather than simply transmitting a function \citep{GaverEtAl2003Ambiguity}. Blythe and colleagues describe anti-solutionist design fiction as a way to resist the assumption that design must always offer direct solutions \citep{BlytheEtAl2016AntiSolutionist}. Lepri, McPherson, and Bowers argue that useless or absurd making can have critical value because it exposes the assumptions that define usefulness \citep{LepriEtAl2020UselessWorthless}. Chindogu, often glossed as the art of the unuseless invention, gives this lineage a particularly apt precedent: the object works, but its working refuses normal utility \citep{PattonBannerot2002Chindogu,Rosenbak2015UselessDesign}. The BRB is not useless, because it records presses. It is not exactly useful either, which is considerate of the theory.

Game culture provides two explicit design influences for this research environment. The Stanley Parable describes itself as a work where ``the rules of how games should work are broken'' and where ``the game plays you'' \citep{StanleyParableOfficial2026}. This matters for BRB because the study similarly stages agency as both real and structured: participants may press or abstain, yet the entire environment has been built to make that freedom legible, tempting, and slightly suspect. Accounting+ describes itself as a ``NIGHTMARE ADVENTURE COMEDY'' for virtual reality \citep{AccountingPlusOfficial2017}, and its relevance lies in the way VR can turn absurd instructions, office-like tasks, bodily presence, and comedic escalation into an interactional system. The BRB adapts these influences to a research setting. It does not ask whether the participant can win. It asks what a participant does when a system gives permission, withholds purpose, counts everything, and maintains a scientifically respectable face.

This humor frame has a methodological boundary. Cultural jokes, memes, and games can motivate the study, but they cannot substitute for evidence about pressing behavior, physiology, attention, or presence. The paper therefore uses humor as an interface concept and design lineage, not as a license to make unsupported empirical claims. The object may be absurd, but the measurements do not get to be.

\subsection{Contribution of the Present Study}

Prior work explains important parts of the phenomenon, but not their convergence in a voluntary red-button encounter. Button history explains how a small gesture can conceal technical and institutional mediation \citep{Plotnick2018}. Nuclear and emergency-control literature explains why the red button is culturally ambivalent, since it can symbolize both catastrophe and interruption \citep{PlotnickWarfare2012,Kwon1996EmergencyStop}. Meme theory explains how a two-button dilemma can circulate as a participatory grammar of conflicted action \citep{Shifman2013Memes,Milner2016WorldMadeMeme}. Humor and absurd-design research explain how expectation violation, uselessness, and ambiguity can become interaction resources rather than defects \citep{Minsky1980Jokes,LepriEtAl2020UselessWorthless}. Embodiment and biofeedback research explain how virtual objects can become more tightly coupled to bodily perception and action \citep{KilteniEtAl2012,LueddeckeFelnhofer2022}. What remains under-specified is how these layers interact when a person is placed in front of a single salient button and told, with complete seriousness, that pressing is entirely optional.

The Big Red Button Institute's VR paradigm addresses that gap through a deliberately minimal and overdetermined encounter. Participants enter a neutral grey VR room containing one large red button positioned at comfortable arm's-length reach. A spoken instruction track preserves agency by stating that participants may press the button, refrain from pressing it, press repeatedly, or never press at all. A visible counter displays cumulative presses, making each action legible as contribution to the study. The physiological manipulation compares live cardiac feedback with sham ECG-like feedback, while respiration modulates ambient lighting. The study then measures pressing behavior, presence, felt peripersonal relation to the button, oneness, and the project-specific construct of secondness. The design is minimal because the room contains very little. It is overdetermined because the object contains too much.

This paper makes three contributions. First, it operationalizes \emph{pressability} as a multilevel construct linking interface history, cultural symbolism, humor, affordance, agency, bodily reach, novelty, and physiological coupling. Second, it documents a VR paradigm for studying voluntary pressing under explicit permission to abstain, which allows non-pressing to count as behavior rather than failure. Third, it frames live-versus-sham cardiac and respiratory feedback as a test of embodied self-adjacency, not as proof that a virtual button becomes part of the participant. The study therefore treats the big red button as an object whose absurdity is experimentally useful. It is ridiculous in the correct direction.

The current manuscript also states its empirical boundary directly. The repository contains a compile-ready protocol and bibliography, but it does not yet contain participant-level datasets, questionnaire exports, raw physiological traces, condition-assignment files, or finalized analysis outputs. The results section therefore remains a truthful scaffold until those materials are archived. This restraint is not a lack of ambition. It is the condition under which a paper about a giant red button can continue to look reviewers in the eye.
